%+TITLE: Rotaciones en 3D
%+DESCRIPTION: Luego de estudiar los grupos, estamos listos para estudiar rotaciones más generales y nuestro primer grupo no abeliano.
\documentclass[11pt,a4paper]{article}

\usepackage[utf8]{inputenc}
\usepackage[T1]{fontenc}
\usepackage[spanish,es-nodecimaldot]{babel}
\usepackage{lmodern}
\usepackage{amsmath,amssymb,mathtools}
\usepackage{graphicx}
\usepackage{xcolor}
\usepackage{geometry}
\usepackage{enumitem}
\usepackage{microtype}

\geometry{margin=2.3cm}
\setlength{\parindent}{0pt}
\setlength{\parskip}{0.65em}
\setlist[itemize]{leftmargin=1.6em,itemsep=0.25em,topsep=0.25em}

\definecolor{title-red}{RGB}{165,25,20}
\definecolor{concept-green}{RGB}{20,125,45}
\newcommand{\concept}[1]{\textcolor{concept-green}{#1}}

\newcommand{\ket}[1]{\left|#1\right\rangle}
\newcommand{\bra}[1]{\left\langle#1\right|}
\newcommand{\braket}[2]{\left\langle #1\middle|#2\right\rangle}
\newcommand{\hc}{\mathrm{h.c.}}
\newcommand{\Tr}{\operatorname{Tr}}
\newcommand{\Diffs}{\mathrm{Diffs}}

\begin{document}

Estudiemos ahora el grupo más sencillo no abeliano y con varios generadores. El grupo de rotaciones en 3 dimensiones.

Una rotación arbitraria la podemos definir como rotar por un ángulo $\alpha$ respecto a un eje $\hat n$, que apunta en una dirección arbitraria. Como $\hat n$ tiene módulo $1$, lo podemos escribir
\[
  \hat n=(\sin\theta\cos\phi,\sin\theta\sin\phi,\cos\theta),
\]
tal que la rotación general depende de tres ángulos: $\alpha,\theta,\phi$.

Cuando $\hat n$ es algún eje coordenado, ya sabemos la respuesta:
\[
  R_{\hat x}(\alpha)=
  \begin{pmatrix}
    1&0&0\\
    0&\cos\alpha&\sin\alpha\\
    0&-\sin\alpha&\cos\alpha
  \end{pmatrix},
  \qquad
  R_{\hat y}(\alpha)=
  \begin{pmatrix}
    \cos\alpha&0&-\sin\alpha\\
    0&1&0\\
    \sin\alpha&0&\cos\alpha
  \end{pmatrix},
\]
\[
  R_{\hat z}(\alpha)=
  \begin{pmatrix}
    \cos\alpha&\sin\alpha&0\\
    -\sin\alpha&\cos\alpha&0\\
    0&0&1
  \end{pmatrix}.
\]

($R_{\hat y}$ tiene esa forma para ser compatible con la regla de la mano derecha.)

Para un eje arbitrario, podemos hacer una composición de rotaciones:

\concept{Primero} rotamos los ejes para que $\hat n$ apunte en la dirección $\hat z$:
\[
  R_{\hat n\to\hat z}=R_{\hat y}(\theta)R_{\hat z}(\phi).
\]
(Tarea: verifique que $R_{\hat n\to\hat z}\hat n=\hat z$.)

\concept{Segundo} rotamos por un ángulo $\alpha$ respecto a $\hat z$: $R_{\hat z}(\alpha)$.

\concept{Tercero} regresamos el vector $\hat n$ a su orientación original: $R_{\hat n\to\hat z}^{-1}$.

Por tanto,
\[
  R_{\hat n}(\alpha)
  =R_{\hat n\to\hat z}^{-1}
   R_{\hat z}(\alpha)
   R_{\hat n\to\hat z}.
\]

Como siempre, es más fácil pensar en términos de invariancias. Las rotaciones dejan invariante el módulo de un vector. Eso es, si $\vec v'=R\vec v$,
\[
  \lVert\vec v'\rVert^2
  =\vec v'{}^T\vec v'
  =(R\vec v)^T(R\vec v)
  =\vec v^{\,T}R^TR\vec v.
\]

Si queremos que $\lVert\vec v'\rVert^2=\lVert\vec v\rVert^2 =\vec v^{\,T}\vec v$, entonces debemos pedir que
\[
  R^TR=\mathbb 1.
\]
Es decir, las matrices ortogonales dejan invariante el módulo de los vectores.

Estas matrices forman un grupo:
\[
  (R_1R_2)^T(R_1R_2)
  =R_2^TR_1^TR_1R_2
  =\mathbb 1,
\]
la identidad es ortogonal, y la inversa de la ortogonal $R^T$ es a su vez ortogonal:
\[
  (R^T)^TR^T=RR^T=\mathbb 1.
\]

Este grupo se llama $O(3)$.

Pero $O(3)$ no contiene sólo rotaciones. Una transformación de paridad también es ortogonal:
\[
  P=
  \begin{pmatrix}
    -1&0&0\\
    0&-1&0\\
    0&0&-1
  \end{pmatrix}.
\]
Pero esta no se puede escribir como una rotación.

De hecho, si tomamos el determinante de $R^TR=\mathbb 1$, obtenemos \[
  \det(R^TR)=1
  \quad\Longrightarrow\quad
  \det R^T\det R=1
  \quad\Longrightarrow\quad
  (\det R)^2=1
  \quad\Longrightarrow\quad
  \det R=\pm1.
\]

Pero una rotación se puede hacer infinitamente cercana a la identidad haciendo $\alpha\to0$. Eso quiere decir que $\det R=1$, mientras que $\det P=-1$.

El subgrupo de $O(3)$ con determinante $1$ se llama $SO(3)$, y corresponde a las rotaciones.

La rotación la podemos escribir en términos de los generadores:
\[
  R_{\hat n}(\alpha)=e^{i\alpha\hat n\cdot\vec J},
\]
donde $\vec J=(J_x,J_y,J_z)$ es un vector de matrices. Para encontrar las $\vec J$, usamos la definición de $SO(3)$ con $\alpha$ infinitesimal:
\begin{align*}
  \mathbb 1
  &=(e^{i\alpha\hat n\cdot\vec J})^T
    (e^{i\alpha\hat n\cdot\vec J})\\
  &\approx(\mathbb 1+i\alpha\hat n\cdot\vec J^{\,T})
    (\mathbb 1+i\alpha\hat n\cdot\vec J)\\
  &\approx\mathbb 1+i\alpha\hat n\cdot
    (\vec J^{\,T}+\vec J).
\end{align*}

Entonces pedimos $\vec J^{\,T}=-\vec J$. Es decir, son matrices antisimétricas. Para que $R$ sea real, $\vec J$ debe ser puramente imaginaria. Una matriz imaginaria antisimétrica tiene $3$ componentes independientes. Es decir, cualquier $\vec J$ se puede expandir en la base
\[
  J_x=
  \begin{pmatrix}
    0&0&0\\
    0&0&-i\\
    0&i&0
  \end{pmatrix},
  \qquad
  J_y=
  \begin{pmatrix}
    0&0&i\\
    0&0&0\\
    -i&0&0
  \end{pmatrix},
  \qquad
  J_z=
  \begin{pmatrix}
    0&-i&0\\
    i&0&0\\
    0&0&0
  \end{pmatrix}.
\]

De hecho estas matrices corresponden a tomar $\alpha$ infinitesimal en las expresiones para $R_{\hat x}$, $R_{\hat y}$ y $R_{\hat z}$.

Con las expresiones para $\vec J$ podemos calcular el álgebra de Lie tomando explícitamente los conmutadores:
\[
  [J_x,J_y]=iJ_z,
  \qquad
  [J_y,J_z]=iJ_x,
  \qquad
  [J_z,J_x]=iJ_y.
\]
Escribiendo $J_1=J_x$, $J_2=J_y$, $J_3=J_z$, podemos resumir esto en
\[
  [J_i,J_j]=i\epsilon_{ijk}J_k.
\]

Resulta que existe otro grupo con la misma álgebra de Lie.

Este grupo es $SU(2)$, el grupo de matrices unitarias $2\times2$,
\[
  U^\dagger U=\mathbb 1,
\]
con determinante $1$,
\[
  \det U=1.
\]

En general, la condición de unitariedad implica
\[
  \det(U^\dagger U)=1,
\]
es decir,
\[
  \det U\det U^\dagger=|\det U|^2=1.
\]

Si yo tengo una matriz unitaria $2\times2$, la puedo escribir
\[
  U=z\widetilde U,
\]
donde $z=\det U$ tal que $|z|^2=1$ y $\det\widetilde U=1$.

Lo que hemos hecho es ``factorizar'' el determinante.

Podemos entonces escribir
\[
  U(2)=U(1)\otimes SU(2),
\]
donde $z$ corresponde a $U(1)$ y $\widetilde U$ a $SU(2)$.

Como ya estudiamos $U(1)$ con las rotaciones 2D, nos enfocamos en $SU(2)$. Contando componentes, vemos que hay tres matrices $SU(2)$ linealmente independientes.

Para encontrar los generadores, hacemos
\[
  U=\mathbb 1+i\alpha G,
  \qquad
  U^\dagger U=\mathbb 1+i\alpha(G^\dagger-G),
  \qquad\Longrightarrow\qquad
  G=G^\dagger.
\]

La condición $\det U=1$ se puede escribir
\begin{align*}
  1=\det U
  &\approx(1+i\alpha G_{11})(1+i\alpha G_{22})
    +\alpha^2G_{12}G_{21}\\
  &\approx1+i\alpha(G_{11}+G_{22})
  \qquad\Longrightarrow\qquad
  \Tr G=0.
\end{align*}

Los generadores son matrices hermíticas sin traza. Hay $3$ independientes:
\[
  \sigma_1=
  \begin{pmatrix}
    0&1\\1&0
  \end{pmatrix},
  \qquad
  \sigma_2=
  \begin{pmatrix}
    0&-i\\i&0
  \end{pmatrix},
  \qquad
  \sigma_3=
  \begin{pmatrix}
    1&0\\0&-1
  \end{pmatrix},
\]
las \concept{matrices de Pauli}.

Éstas satisfacen
\[
  \sigma_i\sigma_j
  =\mathbb 1\,\delta_{ij}+i\epsilon_{ijk}\sigma_k
\]
(verifique haciendo varias multiplicaciones). Resulta que también son unitarias,
\[
  \sigma_i^\dagger\sigma_i=\mathbb 1.
\]
De aquí tenemos que
\[
  [\sigma_i,\sigma_j]=2i\epsilon_{ijk}\sigma_k.
\]

Escribimos el elemento de $SU(2)$ como
\[
  U=\mathbb 1+i\vec\alpha\cdot\vec J,
  \qquad
  J_i=\frac12\sigma_i.
\]

Podemos siempre factorizar números al definir nuestros generadores. De esta manera nos queda
\[
  [J_i,J_j]
  =\frac14[\sigma_i,\sigma_j]
  =\frac12i\epsilon_{ijk}\sigma_k
  =i\epsilon_{ijk}J_k.
\]

Como prometimos, esta es la misma álgebra de Lie de $SO(3)$.

Hay una relación más estrecha entre los dos grupos. En abstracto, podemos escribir un vector como
\[
  \vec v=v_1\hat e_1+v_2\hat e_2+v_3\hat e_3,
  \qquad
  \hat e_i\cdot\hat e_j=\delta_{ij}.
\]

Para la base podemos usar $\hat e_i=\sigma_i$ y para el producto interno
\[
  \vec v\cdot\vec w=\frac12\Tr(vw).
\]
Efectivamente,
\[
  \hat e_i\cdot\hat e_j
  =\frac12\Tr(\sigma_i\sigma_j)
  =\delta_{ij}+i\frac12\epsilon_{ijk}\Tr\sigma_k
  =\delta_{ij}.
\]

Entonces cualquier vector se puede escribir
\[
  \vec v=
  \begin{pmatrix}
    v_z&v_x-iv_y\\
    v_x+iv_y&v_z
  \end{pmatrix}.
\]

Estos vectores se pueden rotar definiendo una rotación respecto al eje $\hat n$ por un ángulo $\alpha$ como
\[
  q=\cos\alpha+i\sin\alpha\,\hat n
  \qquad\text{(escrito como matriz)}.
\]
La rotación resulta ser
\[
  \vec v'=q^{-1}\vec vq.
\]

Haremos un check tomando $\vec v$ en la dirección $\hat x$ y tomando $\hat n=\hat z$.

\begin{align*}
  \vec v'
  &=v_x(\cos\alpha+i\sigma_3\sin\alpha)^\dagger
    \sigma_1(\cos\alpha+i\sigma_3\sin\alpha)\\
  &=v_x\Bigl(
      \cos^2\alpha\,\sigma_1
      +i[\sigma_1,\sigma_3]\sin\alpha\cos\alpha
      +\sigma_3\sigma_1\sigma_3\sin^2\alpha
    \Bigr)\\
  &=v_x\Bigl(
      (\cos^2\alpha-\sin^2\alpha)\sigma_1
      +2\sin\alpha\cos\alpha\,\sigma_2
    \Bigr)\\
  &=\cos2\alpha\,v_x\hat e_1
    +\sin2\alpha\,v_x\hat e_2.
\end{align*}

Nos equivocamos, en realidad la rotación debería ser
\[
  q=\cos(\alpha/2)+i\hat n\sin(\alpha/2).
\]

Esto define un homomorfismo entre $SO(3)$ y $SU(2)$ que mapea dos elementos de $SO(3)$, $\theta$ y $\theta+\pi$, en un elemento de $SU(2)$ dado por $\alpha = 2\theta$.

\end{document}
